\section{三维的规范场：Chern--Simons 作用量}

\subsection{联络与曲率}

设 $P\to M$ 是以紧致单 Lie 群 $G$ 为结构群的主丛，$A$ 是 $P$ 上的联络。在局部平凡化下，
$A=A_i\dd x^i$ 取值为 Lie 代数 $\mathfrak g$ 中。曲率是
\[
F=\dd A+A\wedge A,\qquad F_{ij}=\partial_iA_j-\partial_jA_i+[A_i,A_j].
\]
在规范变换 $A\mapsto A^{\,g}=g^{-1}Ag+g^{-1}\dd g$ 下，曲率协变地变换：$F\mapsto g^{-1}Fg$。
因此 $\Tr(F\wedge\star F)$ 这样的量是规范不变的。

我们还需要固定迹的归一化。取 $G$ 的定义表示，并要求
\[
\Tr(T^aT^b)=\tfrac12\delta^{ab},
\]
即长根平方为 $2$ 的归一化。这个选择决定了下面所有公式里 $k$ 的单位；换表示或换归一化，
$k$ 的整性条件会跟着换。讲义附录里有一张各教材约定对照表，读者在比较不同书时务必先对齐这一项。

\subsection{为什么不用 Yang--Mills 作用量}

最自然的规范场作用量是 Yang--Mills 作用量
\[
L_{\mathrm{YM}}=\int_M\sqrt g\,g^{ik}g^{jl}\Tr(F_{ij}F_{kl})\,\dd^3x.
\]
它需要度规：$\sqrt g$ 是体积元，$g^{ik}g^{jl}$ 把两个曲率指标的上下位置缩并。
这两处对度规的依赖性质不同：前者是体积元，换度规只是改标度，可以用“不积分体积”的方式绕开；
后者把不同分量混合起来，\textbf{绕不开}——换一个度规，作用量里的混合方式就变了。

既然我们要的是一个不依赖度规的配分函数，就不能用 Yang--Mills。需要另一个只依赖微分拓扑的作用量。

\subsection{Chern--Simons 作用量}

在三维流形上，存在一个只用到 $A$ 与 $\dd$ 的三次型：
\begin{equation}
L=\frac{k}{4\pi}\int_M\Tr\left(A\wedge\dd A+\frac23A\wedge A\wedge A\right)
=\frac{k}{8\pi}\int\dd^3x\,\varepsilon^{ijk}\Tr\left(A_i\partial_jA_k+\frac23A_i[A_j,A_k]\right). \tag{2.1}
\end{equation}
这就是 Chern--Simons 作用量。它有三个决定性的性质。

\textbf{第一，它只依赖定向，不依赖度规。} 被积式是三维流形上的 $3$-形式，积分只需要定向，
不需要 $\sqrt g$，也不需要缩并。

\textbf{第二，它的变分给出曲率。} 直接计算给出
\begin{equation}
\delta L=\frac{k}{2\pi}\int_M\Tr(\delta A\wedge F). \tag{2.2}
\end{equation}
因此经典运动方程是 $F=0$：联络必须是平坦的。这个结果的物理含义值得停下来想一想：
三维 Chern--Simons 理论没有传播的自由度，它的经典解是“平坦联络”，
而不是像 Yang--Mills 那样有波。这正是它能成为拓扑场论的原因，也是第三讲相空间为有限维的根源。

\textbf{第三，它在小规范变换下不变，在大规范变换下平移 $2\pi$ 的整数倍。} 对小规范变换
$g$ 接近恒等，$\delta L=0$；对不平凡的大规范变换，$L$ 改变一个常数。这个常数决定了 $k$ 必须取整数。

下面把这三条依次算清楚。

\subsection{变分}

把 (2.1) 代入 (2.2) 并分部积分，可以把 $F$ 提出来。规范场论里常用的恒等式是
\[
\delta\int\Tr(A\wedge\dd A)=\int\Tr(\delta A\wedge\dd A)-\int\Tr(A\wedge\dd\,\delta A)
=2\int\Tr(\delta A\wedge\dd A),
\]
最后一步用了 $\dd A=-A\wedge\dd$（在三维、对一次形式的 $A$）与迹的循环性；
而 $\delta\frac23\int\Tr(A^3)=2\int\Tr(\delta A\wedge A\wedge A)$。
两式相加即得 (2.2)。

\subsection{大规范变换与 $k$ 的整性}

设 $g:M\to G$ 是规范变换，其同伦类由绕数 $m\in\ZZ$ 刻画（因为 $\pi_3(G)=\ZZ$ 对紧致单群成立）。
在 $A\mapsto g^{-1}Ag+g^{-1}\dd g$ 下，经过一段计算，$L$ 的变化为
\begin{equation}
L\longmapsto L+\frac{k}{4\pi}\cdot\frac13\int_M\Tr\left((g^{-1}\dd g)^3\right)
=L+\frac{k}{4\pi}\cdot\frac13\cdot 24\pi^2m
=L+2\pi k\,m, \tag{2.3}
\end{equation}
其中用到了恒等式
\[
m=\frac{1}{24\pi^2}\int_M\Tr\left((g^{-1}\dd g)^3\right)\in\ZZ.
\]
在量子理论里，物理量必须是 $\exp(\ii L)$ 的期望值，而这个量在规范变换下必须不变。
因此要求 $2\pi k m$ 是 $2\pi$ 的整数倍，即
\begin{equation}
k\in\ZZ. \tag{2.4}
\end{equation}
这是 Chern--Simons 理论最著名的性质：耦合常数被量子化。它的根源是 $\pi_3(G)$ 的非平凡性，
与二维共形场论里仿射李代数的中心荷取整数是同一件事的两面。

\subsection{Wilson 线与配分函数}

纽结进入理论的方式是 Wilson 线。给定链环 $C$ 与它的一个表示 $R$，定义
\begin{equation}
W_R(C)=\Tr_R\,P\exp\int_C A_i\dd x^i, \tag{2.5}
\end{equation}
其中 $P$ 表示路径有序。这个量不需要度规，只用到 $A$ 沿曲线的积分，因此是拓扑场论里自然的可观测量。
更一般地，对链环 $C_i$ 与表示 $R_i$，配分函数定义为
\begin{equation}
Z(M;C_i,R_i)=\int\cD A\,\exp(\ii L)\prod_i W_{R_i}(C_i). \tag{2.6}
\end{equation}
这里的路径积分不是严格定义的，但它的半经典展开和手术计算都可以严格化——这是第七讲的内容。

Witten 的中心论断是：当 $G=SU(N)$、所有 $R_i$ 取定义表示时，$S^3$ 上的这个配分函数
恰好给出 Jones 多项式，而两个参数 $N$ 与 $k$ 合起来给出二变量推广\Wp{356}。
当 $G=SO(N)$、$R$ 取 $N$ 维表示时，得到的是 Kauffman 多项式。

\subsection{弱耦合是大 $k$}

最后提醒一个容易混淆的地方。Chern--Simons 理论的“弱耦合极限”是 $k\to\infty$，
不是 $k\to0$。原因是 $k$ 在作用量里是一个整体因子：$k$ 越大，路径积分的相位振荡越剧烈，
半经典近似越好。这与 Yang--Mills 的 $1/g^2$ 展开正好相反，读者在比较不同文献时要特别小心。

\subsection{一个几何解释：Chern--Simons 是四维示性类的边界项}

$k$ 为什么必须是整数，还可以从一个更几何的角度看。在四维流形 $X$ 上，
$p_1=\frac{1}{8\pi^2}\Tr(F\wedge F)$ 是一个整上同调类（由 Atiyah--Singer 指标定理）。
它的一个陈--西蒙斯形式是
\[
\dd\,\omega_{\mathrm{CS}}(A)=\Tr(F\wedge F),
\]
而 (2.1) 的被积式就是 $\omega_{\mathrm{CS}}$ 的积分。如果 $X$ 以 $M$ 为边界，
那么 $\int_X\Tr(F\wedge F)=\int_M\omega_{\mathrm{CS}}$。左边是整数（超荷/瞬子数），
所以右边在 $M$ 上必须整——这正是 $k\in\ZZ$ 的几何来源。
换句话说，Chern--Simons 理论是四维拓扑不变量（第二陈类、瞬子数）在三维的“边界理论”。

这个观点也解释了为什么三维是 Chern--Simons 的“自然维度”：三次型需要一个三维积分，
而它的微分是四维示性类。更高维的推广是 $2n-1$ 维的 Chern--Simons 理论，
但物理上最相关的是 $n=2$。

\subsection{归一化、level 与 $q$}

同一篇论文里的 $k$，在不同的书里可能差一个因子。原因是 Lie 代数的迹有多种归一化：
有人取定义表示、长根平方为 $2$；有人取 Killing 型；有人用“中心荷归一到 $1$”的约定。
本节统一用定义表示、$\Tr(T^aT^b)=\frac12\delta^{ab}$；在这个约定下 (2.1) 的 $k$ 与仿射代数的 level 一致。

量子化之后常用的变量是
\[
q=\exp\left(\frac{2\pi\ii}{k+h^\vee}\right),
\]
其中 $h^\vee$ 是对偶 Coxeter 数（$SU(N)$ 时为 $N$）。$q$ 的单位根性质是“可积表示只有有限多个”
的直接体现：当 $q$ 是单位根时，量子群的表示范畴退化，只剩下有限多个可积对象。
Jones 多项式里的 $t$ 与这个 $q$ 的关系是 $t=q^{2}$（不同文献的约定略有差别），
读者在比较时必须先对齐 $k$、$h^\vee$ 与 $q$ 的定义。

\subsection{路径积分可算吗}

(2.6) 的路径积分不是严格定义的对象：三维流形上的规范场测度没有已知的构造。
但 Chern--Simons 理论有一个特殊性质：它的路径积分只依赖拓扑，
因此在半经典极限（$k\to\infty$）与手术分解下都可以“算”。Witten 的做法是：
把三维流形沿曲面切开，用曲面上的态空间做粘合；曲面上的态空间是有限维的，
所以粘合是一个有限维的线性代数操作。这样，路径积分就化归为一组自洽的有限维数据
（$S$ 矩阵、融合系数、标架相位），而这些数据由二维共形场论提供。

换句话说，这个理论不是靠“把路径积分算出来”来定义的，而是靠一组公理（TQFT 公理）
加上一组由共形场论给出的解来定义的。第七讲会把这组数据写具体。

\subsection{带边界的情形与边界项}

前面都假设 $M$ 无边。一旦 $M$ 有边界 $\partial M=\Sigma$，变分 (2.2) 会多出一个边界项：
\begin{equation}
\delta L=\frac{k}{2\pi}\int_M\Tr(\delta A\wedge F)
-\frac{k}{4\pi}\int_\Sigma\Tr(\delta A\wedge A). \tag{2.7}
\end{equation}
要让变分良定义，必须指定边界上 $A$ 的什么分量固定。常见的选择是固定 $A$ 在 $\Sigma$ 上的
一个复结构分量（即固定 $A^{(0,1)}$）；这时边界项消失，体内方程仍是 $F=0$。

这个边界项在故事里很重要。第五讲把纽结写成穿刺点时，穿刺点附近的小圆盘就是带边界的曲面；
边界项决定了“在边界上插入什么态”。第七讲的手术公式，本质上也是把边界上的态求和掉。
换句话说，Chern--Simons 理论不是“只活在三维内部”的理论：它的边界携带一个二维的态空间，
这正是它能与共形场论对接的技术原因。

\subsection{$U(1)$ 情形：一个可以完全算清的例子}

把群换成 $U(1)$ 时，一切都可以显式写出。$A$ 是一个普通一形式，$F=\dd A$，作用量是
\[
L=\frac{k}{4\pi}\int_M A\wedge\dd A .
\]
运动方程 $\dd A=0$，即 $A$ 是闭形式。把 $A$ 分解成“恰当部分 + 调和部分”，
恰当部分可以规范掉，只剩下调和部分；于是理论化归为有限维的辛向量空间上的高斯积分。
在曲面 $\Sigma$ 上，$U(1)$ 理论的 Hilbert 空间是 theta 函数空间，维数为 $k^g$；
环面上（$g=1$）维数为 $k$，与第四讲的截面计数一致。

$U(1)$ 情形里 $k$ 的整性同样来自大规范变换的绕数，与第二讲完全平行；
但这里没有非阿贝尔的复杂性，所以可以拿来检验每一个公式。读者若对某一步不放心，
先退到 $U(1)$ 算一遍，通常能看清结构。

\subsection{微扰展开：$1/k$ 展开给出什么}

$k\to\infty$ 时，路径积分可以用驻点法展开，展开参数是 $1/k$，展开的图形是三维费曼图。
有一个漂亮的结果：这些费曼图给出的不变量，恰好是 Vassiliev 不变量（有限型不变量）的各级；
换句话说，Chern--Simons 微扰论给出纽结不变量的一个无穷族，Jones 多项式只是其中最粗的一层。

这件事有两个意义。第一，$1/k$ 展开不是形式幂级数，而是有几何意义的：
每一项对应一类奇点的消解。第二，它说明为什么 Jones 多项式不够：
有限型不变量的层级比单个多项式丰富得多。

不过微扰展开有一个技术困难：三维费曼图会出现紫外发散，需要正则化；
而正则化会引入标架依赖，这正是第七讲标架反常的来源。也就是说，
微扰论的自洽性要求标架结构，而标架在拓扑场论里本来就是一个必须处理的数据。

\subsection{为什么三维 Chern--Simons 是可解的}

一个三维量子场论通常无法精确求解：自由度太多，相互作用太强。Chern--Simons 之所以例外，
原因可以一句话概括：它没有局域自由度。

作用量 (2.1) 不含度规，也不含 $F$ 的平方项，所以它不是一个“传播”的理论：
没有引力子、没有规范玻色子的传播子，唯一的动力学是整体（拓扑）的。
把流形沿曲面切开后，每一侧的态空间是有限维的；所有的“物理”都集中在边界上。
这就是为什么它可以用二维共形场论的数据完全描述。

这个观察在技术上体现为两点。第一，相空间是有限维的（第三讲）：它不是“场的空间”，
而是“解的空间”。第二，配分函数只依赖流形的拓扑（手术不变量，第七讲）。
物理上，这类理论叫拓扑场论（TQFT）；Witten 1988 年的另一篇文章给出了 TQFT 的公理，
而本文是它最成功的例子。

从另一个角度看，Chern--Simons 与二维模型的关系是“同一组数据的两种组织方式”：
三维侧是拓扑的、整体的，二维侧是共形的、局域的。这个对偶关系不是本文发明的，
但本文第一次说明它为什么成立：三维理论在曲面上的态空间，就是二维理论的共形块。

\subsection*{练习}
\begin{enumerate}[leftmargin=*]
  \item 验证 (2.2)：从 (2.1) 出发分部积分，注意三维中 $A\wedge\dd A$ 与 $A\wedge A\wedge A$ 的系数 $\frac23$ 如何分配。
  \item 证明 $m=\frac{1}{24\pi^2}\int\Tr((g^{-1}\dd g)^3)$ 对 $g:S^3\to SU(2)$ 取整数，并计算 $g$ 为恒等映射的倍数时的值。
  \item 对 $G=SU(2)$，用 $\Tr(T^aT^b)=\frac12\delta^{ab}$ 写出作用量在 $A=A^aT^a$ 下的显式形式，并检验它与 (2.1) 的系数一致。
  \item 说明为什么 $k\to0$ 不是弱耦合极限。
\end{enumerate}

\begin{quote}
\small
\textbf{本讲要点。} Chern--Simons 作用量只依赖定向；(2.2) 给出 $F=0$ 的经典解；
大规范变换 (2.3) 强制 $k\in\ZZ$；Wilson 线 (2.5) 与配分函数 (2.6) 是纽结不变量的候选；
$k$ 的几何来源是四维示性类的边界项。
\end{quote}
